\section{A Cost Model for Salt Cardinality}
\label{sec:costmodel}

\subsection{Formulation}

Under selective salting with cardinality $k$, total stage time decomposes into
three parts.

\emph{The critical task.} The hot key's $H$ rows are spread over $k$
sub-partitions, so the heaviest task processes approximately $H/k$ rows. Write
$a$ for the marginal cost of a row on the critical path.

\emph{Replication.} The $P$ matching probe-side rows are emitted $k$ times
instead of once, adding $P(k-1)$ rows to shuffle, sort and merge. Write $b$ for
the marginal cost of a replicated row.

\emph{Everything else.} Scanning both inputs, joining the cold keys, and the
terminal aggregation are all independent of $k$. Collect them into $c$.

\begin{equation}
T(k) \;=\; a\,\frac{H}{k} \;+\; b\,P\,(k-1) \;+\; c ,
\qquad k \ge 1 .
\label{eq:model}
\end{equation}

\subsection{The optimum}

$T$ is the sum of a convex decreasing term and a linear increasing term, hence
strictly convex on $k>0$, with a unique interior minimum. Differentiating,

\begin{equation}
\frac{dT}{dk} \;=\; -\frac{aH}{k^{2}} \;+\; bP \;=\; 0
\qquad\Longrightarrow\qquad
\boxed{\;k^{*} \;=\; \sqrt{\frac{aH}{bP}} \;=\; \sqrt{\gamma\,\frac{H}{P}}\;}
\label{eq:kstar}
\end{equation}

where $\gamma = a/b$ is the ratio of the two marginal costs. Four properties of
Equation~\eqref{eq:kstar} are worth drawing out.

First, \textbf{it depends on the workload only through $H/P$}, the ratio of
hot-key rows on the two sides. That ratio is cheap to estimate and is exactly
what a frequency sketch on the join key already yields.

Second, \textbf{$\gamma$ ought to be a property of the engine and the hardware
rather than of the query}: row-processing cost on the critical path and
replicated-row shuffle cost both scale with the machine, so their ratio should
be more stable than either alone, and one calibration should serve many
workloads. We state this as a hypothesis rather than a result because our own
measurements do not support it --- see Section~\ref{sec:results}, where fitted
$\gamma$ varies by \GammaSpread$\times$ across three workloads on one machine.
That instability is diagnostic of the fits being dominated by the $k=1$ to
$k=2$ step rather than by the replication term, and it is the second reason,
alongside the unresolved convexity, that we do not claim the closed form is
validated.

Third, \textbf{$k^{*}$ falls as $1/\sqrt{P}$}. Thickening the probe side's hot
key makes every additional salt value more expensive while leaving straggler
relief unchanged, so the optimum retreats. This is the model's sharpest
falsifiable prediction, and Section~\ref{sec:results} tests it directly by
sweeping $\mu$ over a factor of 32.

Fourth, and least intuitively, \textbf{$k^{*}$ grows only as the square root of
the skew}. Doubling the hot key's dominance does not double the salt
cardinality; it multiplies it by $\sqrt{2}$. This is the structural reason
aggressive salting is so often counterproductive: intuition says respond to
severe skew with large $k$, and the mathematics says respond with a modestly
larger one.

\subsection{Two ceilings}

Equation~\eqref{eq:kstar} is unconstrained, and taken literally it will
sometimes recommend a $k$ that buys nothing at all.

\subsubsection{Saturation}
Once the hot partition has been divided below the size of an ordinary
partition, the hot task stops being the straggler --- some other partition now
finishes last --- and further division relieves nothing while continuing to pay
replication. From Equation~\eqref{eq:rho}, that point is
\begin{equation}
k_{\text{sat}} \;=\; \rho \;=\; \frac{H}{N/p} .
\end{equation}
\textbf{Salting beyond $\rho$ is always waste.} The bound is immediate from the
model, requires no calibration, and we have not found it stated in the
practitioner literature.

\subsubsection{Parallelism}
Sub-partitions beyond the number of concurrently schedulable task slots $C$
queue rather than overlap, so the straggler term stops falling at $k=C$ while
the replication term keeps rising:
\begin{equation}
k_{\text{par}} \;=\; C .
\end{equation}

\subsubsection{The operational rule}
\begin{equation}
k_{\text{eff}} \;=\; \max\Bigl(1,\;\bigl\lfloor \min(k^{*},\,\rho,\,C) \bigr\rceil\Bigr) .
\label{eq:keff}
\end{equation}

Which of the three binds is itself diagnostic. If $k^{*}$ binds, the workload
is replication-limited and the probe side is the thing to shrink. If $\rho$
binds, the skew is mild relative to the partition count and raising $p$ would
serve better than salting. If $C$ binds, the cluster is the limit and salting
further is pure overhead --- the common case on small clusters, and the one in
which tutorial defaults do the most damage.

\subsection{Uniform salting}

For uniform salting the replication term scales with the whole probe table, so
Equation~\eqref{eq:model} becomes $T_{u}(k) = aH/k + bM(k-1) + c$ and
$k_{u}^{*} = \sqrt{\gamma H/M}$. Since $M \gg P$, the optimum is
correspondingly smaller --- often below $2$, meaning uniform salting is
frequently not worth doing at all. This falls out of the model rather than
having to be discovered empirically, and it is a good illustration of why
pricing the technique is worth the trouble.

\subsection{Calibration}

Given measurements $\{(k_i, T_i)\}$, Equation~\eqref{eq:model} is \emph{linear}
in $(a,b,c)$ under the basis $[H/k,\; P(k-1),\; 1]$. Fitting is therefore
ordinary linear least squares: no initial guess, no local minima, no optimiser
hyperparameters, and a fit that is a deterministic function of the
measurements. Three distinct values of $k$ suffice to identify three
coefficients; we use seven.
